
\sourcenote{Genre: core rule (tier 1) --- an ordinary mathematical proof
at Putnam-exam-solution detail. The rule's mechanization is
\emph{conditional}: the section descends the invariance obligation from
an opaque interface fact to three named hypothesis families ---
integrated facet-row clauses, grid geometry of the block, arena
invariance --- and the ledger remark at the end states exactly what
remains and why. Every hypothesis is used and remarked on; the
formal-anchor table maps each step to the Lean declaration that carries
it.}

\textbf{Setting and claim.}\enspace Fix a forward semiflow $\varphi$ on
a state space $X$ ($\Phi_0 = \mathrm{id}$,
$\Phi_{s+t} = \Phi_t \circ \Phi_s$; the \rulename{QBasin} section's
setting, which also fixes the $\mathsf{QInvScope}$ denotation: $S$
forward-invariant and $X_0 \subseteq S$). The rule's data: a
\emph{block} $N$ --- the realization of a box-union over the
certificate's declared rational grid, a compact closed set --- sitting
in an \emph{arena} $D \supseteq N$; per internal facet of $N$, one
strictly-inward sign row of the base calculus (a
\rulename{BernCone}-family sequent asserting the field's normal
component points into the block, with an explicit positive rational
margin; facets on $\partial D$ owe nothing); and one containment row
$X_0 \subseteq N$. The claim: these data warrant
$\mathsf{QInvScope}(N, X_0)$.

\begin{blemma}[rule-level assembly]\label{lem:qbi-assembly}
If $N$ is forward-invariant and $X_0 \subseteq N$, then
$\mathsf{QInvScope}(N, X_0)$ holds.
\end{blemma}
\begin{proof}
The judgement \emph{is} that conjunction. The lemma is stated to fix
where the rule-level content ends: everything below descends the
invariance conjunct toward the certificate's rows.
\end{proof}

\begin{blemma}[the first-exit skeleton: closed $+$ continuous $+$
boundary dwell $\Rightarrow$ invariant]\label{lem:qbi-dwell}
Suppose $N$ is closed, every trajectory $t \mapsto \Phi_t x$ is
continuous, and every boundary point dwells: for each
$x \in N \cap \partial N$ there is $\delta > 0$ with
$\Phi_s x \in N$ for all $0 \leq s \leq \delta$. Then $N$ is
forward-invariant.
\end{blemma}
\begin{proof}
First, \emph{every} point of $N$ dwells: an interior point $x$ has
$\Phi_0 x = x \in \mathrm{int}\,N$, so continuity of the trajectory at
time $0$ confines a short prefix to $\mathrm{int}\,N$ --- the dwell
hypothesis is needed only at the frontier, which is where the
certificate's rows live.

Now fix $x \in N$ and $t \geq 0$, and suppose toward a contradiction
that $\Phi_t x \notin N$. Consider the set of \emph{prefix-confined
times}
\[
S \;:=\; \{\, u \geq 0 \ :\ \Phi_s x \in N \ \text{for all } s \leq u
\,\}.
\]
$S$ contains $0$, and every element of $S$ is $\leq t$ (an element
$> t$ would confine time $t$ itself). Let $T := \sup S \leq t$. Every
time $s < T$ is confined (it is below some element of $S$). The sup
itself is confined: for $T > 0$ the preimage of $N$ under the
continuous trajectory is closed and contains $[0, T)$, hence contains
$T \in \mathrm{cl}\,[0,T)$; for $T = 0$ it is the hypothesis
$x \in N$. So $T \in S$. But now the dwell at the point $\Phi_T x \in N$
supplies $\delta > 0$ with $\Phi_s(\Phi_T x) \in N$ for
$0 \leq s \leq \delta$,
i.e. (semiflow law) $\Phi_u x \in N$ for $T \leq u \leq T + \delta$;
combined with $T \in S$ this puts $T + \delta \in S$, contradicting
$T = \sup S$ with $\delta > 0$.
\end{proof}

\begin{bremark}[each skeleton hypothesis is essential]
\emph{Closedness}: on $X = \R$ with $\Phi_t x = x + t$ and the
half-open $N = [0, 1)$, every point of $N$ dwells (interior points by
continuity; the one frontier point of $N$, namely $0$, dwells since
$\Phi_s 0 = s \in [0,1)$ for small $s$ --- the other frontier point
$1$ is not in $N$), yet the trajectory from $0$ exits: the proof fails
exactly at
``the sup itself is confined'', which is where closedness is consumed.
\emph{Trajectory continuity}: with a discontinuous time-dependence the
preimage argument dies and a trajectory may jump across the frontier
without ever presenting a boundary point to the dwell hypothesis; in
this stratum continuity is a model-class fact (polynomial-drift flows),
not a theorem. \emph{Dwell at the frontier}: dropping it re-opens exit
through the boundary --- it is precisely the obligation the facet rows
discharge (next lemma).
\end{bremark}

\begin{blemma}[the facet-row bridge: certified rows $\Rightarrow$
boundary dwell]\label{lem:qbi-bridge}
Suppose the certificate supplies, for block $N$ in arena $D$: finitely
many \emph{facet observables} $\ell_1, \dots, \ell_n : X \to \R$ (the
affine facet polynomials), margins $\varepsilon_i > 0$, and certified
regions $F_i$ (the facet boxes, thickened to full dimension), such that

\emph{(row, integrated form)} along any trajectory segment confined to
$F_i$, $\ell_i$ grows at rate at least $\varepsilon_i$:
if $\Phi_s x \in F_i$ for all $s \leq t$, then
$\ell_i(\Phi_t x) \geq \ell_i(x) + \varepsilon_i\, t$; and

\emph{(local wedge test)} every frontier point $x$ of $N$ has an
\emph{active set} $A \subseteq \{1, \dots, n\}$ --- each active
observable vanishing at $x$, each active $F_i$ a neighbourhood of $x$
--- and a neighbourhood $U$ of $x$ such that every $y \in U$ with
$y \in D$ and $\ell_i(y) \geq 0$ for all $i \in A$ lies in $N$.

If moreover $N \subseteq D$, $D$ is forward-invariant, and trajectories
are continuous, then every boundary point of $N$ dwells.
\end{blemma}
\begin{proof}
Fix $x \in N \cap \partial N$ with active set $A$ and wedge
neighbourhood $U$. The set $U \cap \bigcap_{i \in A} F_i$ is a finite
intersection of neighbourhoods of $x$, hence a neighbourhood; continuity
of the trajectory at time $0$ (where it sits at $x$) confines a prefix
$[0, \delta]$ to it, some $\delta > 0$. Now take any
$0 \leq s \leq \delta$
and check the wedge test's three conditions at $y = \Phi_s x$:
membership in $U$ is the confinement; membership in $D$ is arena
invariance applied to $x \in N \subseteq D$; and for each active $i$,
the whole prefix up to $s$ is confined to $F_i$, so the integrated row
gives $\ell_i(\Phi_s x) \geq \ell_i(x) + \varepsilon_i s
= 0 + \varepsilon_i s \geq 0$, using exactly the activation fact
$\ell_i(x) = 0$. The wedge test re-admits $\Phi_s x \in N$. Inactive
facets never appear: the test quotes only active ones.
\end{proof}

\begin{bthm}[soundness of \rulename{QBlockInv}, at the discharged
floor]\label{thm:qbi-sound}
Suppose: $N$ is closed with $N \subseteq D$; $D$ is forward-invariant;
trajectories are continuous; the certificate's facet rows satisfy
Lemma~\ref{lem:qbi-bridge}'s two clauses; and $X_0 \subseteq N$. Then
$\mathsf{QInvScope}(N, X_0)$ holds.
\end{bthm}
\begin{proof}
Lemma~\ref{lem:qbi-bridge} delivers the boundary dwell;
Lemma~\ref{lem:qbi-dwell} converts it to forward invariance of $N$;
Lemma~\ref{lem:qbi-assembly} assembles the judgement with the
containment row.
\end{proof}

\begin{bremark}[the wedge test is one-directional by design]
\label{rem:qbi-wedge}
At a \emph{reentrant} corner of a box-union block --- an inner corner
of an L-shape --- the block is locally the \emph{union} of two
half-spaces $\{\ell_1 \geq 0\} \cup \{\ell_2 \geq 0\}$, while the
wedge test's hypothesis set is their \emph{intersection}: the
implication ``both $\ell_i \geq 0$ $\Rightarrow$ $y \in N$'' holds
(the intersection sits inside the union), but its converse is false
there --- points of $N$ with one active observable negative exist on
the other branch. The biconditional form of the test would therefore
be unsatisfiable at reentrant corners, while the one-directional form
is exactly what the first-exit argument consumes
(Lemma~\ref{lem:qbi-bridge} never needs to \emph{characterize} the
frontier, only to re-admit prefix points). A certificate whose rows
are jointly unsatisfiable at a boundary point --- checkerboard blocks
touching at a corner and demanding opposite strict signs there ---
simply makes the bridge's hypotheses unfulfillable: soundness only
consumes what is given.
\end{bremark}

\begin{bremark}[$\partial D$-faces owe nothing; the arena regress
terminates]
Facets of $N$ on the arena boundary carry no rows: in
Lemma~\ref{lem:qbi-bridge} they enter only through the conjunct
$y \in D$, which arena invariance supplies. The arena's own invariance
is either a nested $\mathsf{QInvScope}$ citation (a chart sub-block
certified by its own facet rows --- this same rule, one level up) or
the model-class fact that a chart's denotation confines trajectories to
its domain; the regress terminates at the whole space, which is
invariant vacuously.
\end{bremark}

\begin{bremark}[strictness of the margins: where it is and is not
consumed]\label{rem:qbi-strict}
The invariance conclusion of this section survives
$\varepsilon_i = 0$: the integrated clause then says $\ell_i$ is
non-decreasing along confined segments, and growth-from-zero still
yields $\ell_i \geq 0$ --- consistent with the classical invariance
theorem (Nagumo \cite{nagumo1942}), which needs only subtangentiality.
The explicit positive rational margin is nonetheless mandatory in the
certificate convention, for two reasons that live outside this
section. First, exactness: a \emph{strict} field sign cannot be
certified directly by the non-strict enclosure predicates of the base
calculus; ``non-strict row on $\mathcal{A}\ell - \varepsilon$ with
explicit $\varepsilon > 0$'' is how strictness becomes checkable in
exact $\Q$. Second, the same rows are consumed by the Morse-family
rules, where strictness is genuinely essential --- no orbit rides a
facet, distinct fibres' invariant parts are disjoint, and the facet
margin bounds admissible chain jumps (``strictness pays a second
time''; see the \rulename{QMorseDischarge} section's ledger). A margin
quietly dropped to zero would leave this rule's conclusion intact and
break those.
\end{bremark}

\begin{bremark}[the hypothesis ledger --- what remains, named]
\label{rem:qbi-ledger}
The chain above discharges the rule's invariance obligation down to
three hypothesis families, and the rule's Lean column renders with the
hypothesis-flavoured check because of exactly these:

\emph{(a) the integrated row clause.} The certificate's rows are
pointwise: per facet box, a base-calculus sequent bounding the field's
certified rate, $\mathcal{A}\ell_i \geq \varepsilon_i$ (the generator
$\mathcal{A} = f \cdot \nabla$ of the deterministic model class).
Lemma~\ref{lem:qbi-bridge} consumes the \emph{integrated} clause along
trajectory segments. The bridge between them is one-variable calculus
--- a mean-value argument along the trajectory --- whose surviving
analytic content is the identification of the certified polynomial
rate with the trajectory-wise derivative on the region: a model-class
profile warrant (for smooth flows the identification is literal;
switched and hybrid profiles replace it by one-sided derivative
bounds, which is their own seam, deliberately not this section's).

\emph{(b) grid geometry of the block.} Closedness of $N$,
$N \subseteq D$, and the wedge test are finite geometry of the
box-union over its declared rational grid --- recomputable
combinatorics, and proved outright for grid blocks in the accompanying
development's grid-geometry module (not bound by this section, whose
theorems take (b) as hypotheses at the interface).

\emph{(c) arena invariance.} As in the regress remark: nested
citation, or model-class warrant, terminating at the whole space.

None of the three is an opaque dynamical assumption; each names a
specific fact with a specific owner --- (a) the profile stratum, (b)
the checker's recompute (and the grid module), (c) the arena's own
certificate. That is the sense in which this core rule is
\emph{conditionally} mechanized: the cited theorems are sorry-free,
and their hypotheses are the ledger above.
\end{bremark}

\begin{bremark}[relation to the literature]
The block-with-inward-facets certificate is the exact-arithmetic
descendant of Conley theory's isolating blocks
\cite{conleyeaston1971}: compact sets whose boundary behaviour is
transverse, so that dynamical claims reduce to boundary sign checks.
The invariance argument itself is Nagumo's theorem
\cite{nagumo1942} in first-exit form, specialized to polyhedral sets,
where the subtangentiality condition becomes finitely many polynomial
sign rows --- the form set-invariance control has long exploited
\cite{blanchini1999}. What the certificate stratum adds is the
discipline: strict signs via non-strict predicates with exact rational
margins (Remark~\ref{rem:qbi-strict}), and a one-directional wedge
test chosen so that box-union (nonconvex) blocks are admissible
without any frontier characterization
(Remark~\ref{rem:qbi-wedge}).
\end{bremark}

\begin{anchors}
Lemma \ref{lem:qbi-assembly} & \texttt{qBlockInv\_sound} (QtslSoundness.lean)\\
Lemma \ref{lem:qbi-dwell}, first-exit skeleton & \texttt{invariant\_of\_boundaryDwell} (Discharge.lean); composed as \texttt{qBlockInv\_sound\_of\_dwell}\\
Lemma \ref{lem:qbi-bridge}, row shadow + wedge test & \texttt{HasFacetRows} (the definition); \texttt{HasFacetRows.boundaryDwell}\\
Theorem \ref{thm:qbi-sound}, end-to-end & \texttt{qBlockInv\_sound\_of\_facetRows}\\
arena regress terminus & \texttt{invariant\_univ}; nested-citation path \texttt{invariant\_of\_invScopeSat} (development)\\
ledger item (a), the mean-value bridge & \texttt{integrated\_ge\_of\_flowDeriv\_ge} + \texttt{qBlockInv\_sound\_of\_rows} (development, not bound here)\\
ledger item (b), grid facts & \texttt{Grid.wedge\_containment}, \texttt{Grid.qBlockInv\_sound\_of\_grid} (development, not bound here)\\
facet rows & strictly-inward \rulename{BernCone}-family sequents with explicit rational margin (base kernel)\\
\end{anchors}
